\section{Method}
\label{sec:method}

\subsection{Preliminary: Qwen-style SDPA-output gating}
\label{subsec:qwen-gate}

A standard attention block produces queries, keys, values
$Q,K,V = XW_Q, XW_K, XW_V$ and computes the scaled dot-product
attention output
\begin{equation}
  \text{Attn}(Q,K,V) = \text{softmax}\!\left(\tfrac{QK^\top}{\sqrt{d_k}}\right) V.
  \label{eq:sdpa}
\end{equation}
After concatenation over heads, the result is linearly projected by
$W_O$ to give the block output. \citet{qiu2025gated} insert a gate
\emph{between} the SDPA output and $W_O$:
\begin{equation}
  Y' = Y \odot \sigma(X W_\theta),
  \label{eq:qwen-gate}
\end{equation}
where $Y$ is the concatenated multi-head SDPA output, $X$ is the
pre-normalised hidden state, $\sigma$ is the element-wise sigmoid, and
$W_\theta \in \mathbb{R}^{d_\text{model} \times (q \cdot d_k)}$
parameterises the gate. The gate is \emph{head-specific} (different
gating vector per head), \emph{elementwise} (per-dimension inside each
head), and \emph{multiplicative}. Qiu et al.~sweep the gate through
five candidate positions ($G_1$ at the SDPA output, $G_2$ at the value
projection, $G_3$ at the key, $G_4$ at the query, $G_5$ at the dense
output) and find $G_1$ dominant on every benchmark tested: PPL,
HellaSwag, MMLU, GSM8k, C-eval.

\subsection{Metadata-aware gate}
\label{subsec:metadata-gate}

Our modification of \eqref{eq:qwen-gate} is to extend the gate's input
with learner-metadata embeddings, broadcast across the sequence:
\begin{equation}
  Y' = Y \odot \sigma\!\bigl(\,[\,X \,;\, e_\text{CEFR}\,;\, e_\text{L1}\,]\, W_\theta\,\bigr),
  \label{eq:metadata-gate}
\end{equation}
with
\begin{itemize}
  \item $e_\text{CEFR} \in \mathbb{R}^{d_c}$ a learned embedding over
    $\{A1, A2, B1, B2, C1, C2, \texttt{unk}\}$, $d_c = 16$,
  \item $e_\text{L1} \in \mathbb{R}^{d_l}$ a learned embedding over the
    EFCAMDAT L1 inventory augmented with an \texttt{unk} slot,
    $d_l = 32$,
  \item both embeddings tiled across sequence length $n$ before
    concatenation along the feature axis with $X$.
\end{itemize}
The gate parameter tensor grows to
$W_\theta \in \mathbb{R}^{(d_\text{model} + d_c + d_l) \times (q \cdot d_k)}$.
For GPT-2 Base ($d_\text{model}{=}768$, $q{=}12$, $d_k{=}64$) this adds
roughly $(48) \times (768) \approx 37\text{k}$ parameters per layer of
gate, i.e.~0.4M over twelve layers, on top of the 2M added by a vanilla
Qwen-style gate. Total overhead remains under 2\% of GPT-2 Base's 117M
parameters.

\paragraph{Initialisation.}
Following \citet{qiu2025gated} we zero-initialise $W_\theta$ so that
$\sigma(\cdot) = 0.5$ at step 0; for stability we actually initialise
the final linear layer with zeros and the bias with a large positive
constant so that $\sigma(\cdot) \approx 1$ at step 0, making the gated
model equivalent to the ungated base model before any gradient updates.
This is important for continual pre-training: the learner-tuned loss
should not regress the base model in the first few steps.

\paragraph{Missing-metadata inference.}
At eval time on corpora without CEFR or L1 labels (or with
out-of-inventory L1), the corresponding embedding is routed to the
\texttt{unk} slot. This is a deliberate design choice: it lets the
model be deployed on texts whose metadata is absent, at the cost of
reducing the gate to a learned average over the training distribution.
For genuine zero-shot transfer to unseen L1s we recommend instead v07
(typological features, Section~\ref{sec:experiments}).

\paragraph{Why $G_1$?}
Qwen's explanation for why $G_1$ dominates $G_2$--$G_5$ is that the
gate interposes a non-linearity between the two otherwise-linear
projections $W_V$ and $W_O$ (Eqs. 7--8 of their paper). For L2 NWP this
matters mechanistically: the non-linearity is what lets the metadata
signal re-weight the value-aggregated context \emph{before} it is mapped
back into the residual stream, i.e.~at the last point where contextual
information is still disentangled per head. Gating earlier (at $Q$, $K$,
or $V$) modulates how information is \emph{gathered}; gating at $G_1$
modulates what is \emph{retained}, which is the operation a
proficiency-aware learner model plausibly needs.
